
\paragraph{Evidence and purpose.}
These are complete written arguments, not Lean-checked or externally
peer-reviewed results. No novelty claim is made. The result justifies a
finite approximation route for the actual unrestricted minimax problem. It
does not identify its rate or make a finite computation an exact certificate
for an unbounded experiment.

\paragraph{Model and computational access.}
Fix a known probability law $\mu$ on $(0,1/2]$ and known $0<\rho\le1/4$.
A fresh coin draws $D\sim\mu$ and $S\mid D=d\sim\mathrm{Ber}(q(d))$, and
supplies conditionally iid $\mathrm{Ber}(1/2+(2S-1)D)$ bits. The allocation
$q$ is measurable, $0\le q\le1$, and $\theta(q)=\int q\,d\mu\le\rho$.
An algorithm sees bits and identities only, and may create, interleave,
revisit, abandon and stop adaptively, with arbitrary independent seeds.
Every actual bit is charged; deterministic computation is not charged.

For the mathematical approximation we grant exact values of the masses of
the finitely many chosen bins, and exact real arithmetic on the resulting
finite coefficients. This is an explicit known-law oracle convention.
An arbitrary measure given without an effective representation need not have
computable bin masses. An executable certified computation would need exact
coefficients or separately certified coefficient errors; neither effective
computability for every measure nor a finite-program solve is claimed here.
The quadratic law with rational endpoints has exact rational bin masses.
Knowing the law never reveals an individual coin's latent magnitude.

Write $V_\mu(B)=R_{\mu,\rho}(B)^2$ for minimax squared error at uniform
expected budget $B>0$. We may restrict all outputs to $[0,\rho]$: projection
onto this interval decreases squared loss for every admissible target and
does not change work. Let $V_{\mu,N}(B)$ impose also at most $N$ total flips,
where $N\ge1$ is an integer. This is a total cap, not a per-coin cap.

Choose a finite interval partition of $(0,1/2]$, retain its positive-mass bins
$A_1,\ldots,A_m$, and choose pairwise-distinct representatives
$d_j\in(0,1/2]$ such that
$|d-d_j|\le\eta$ on $A_j$. Put $\pi_j=\mu(A_j)$ and
$\mu_\eta=\sum_j\pi_j\delta_{d_j}$. Such partitions and distinct interior
representatives exist for every $\eta>0$; discarded bins are $\mu$-null and
all allocation lifts may be defined arbitrarily on them. For
$0<\gamma\le1$, define the finite parameter set
\[
 Q_\gamma=\{q\in(\gamma\mathbb Z_{\ge0})^m:
                    0\le q_j\le1,\ \sum_j\pi_jq_j\le\rho\}.
\]
Let $G_{N,\eta,\gamma}(B)$ be the minimax squared error on this finite
parameter set and rounded magnitude law, using the full adaptive population
experiment with at most $N$ total flips and uniform expected work at most
$B$. The actual representatives and bin masses, not merely $\eta$, are part
of this notation. The finite experiment still hides every sampled type.

\begin{theorem}[Explicit expected-budget approximation]\label{thm:approx}
Put
\[
 C=N^2(\eta+\gamma),\qquad
 E=\rho^2N(\eta+\gamma)+2\rho\gamma.
\]
Then
\[
 \left|V_\mu(B)-G_{N,\eta,\gamma}(B)\right|
 \le \frac{\rho^2B}{N}+E+
                  \frac{\rho^2 C}{4(B+C)}.                 \tag{1}
\]
In particular, for fixed $\mu,\rho,B>0$, taking $N\to\infty$ and partitions
with $\eta=N^{-3}$ and grids $\gamma=N^{-3}$ makes the right side tend to
zero. Each approximating experiment has a finite-dimensional convex
minimax formulation retaining all adaptive coin allocations and exactly
the specified uniform expected-work constraint.
\end{theorem}

\begin{lemma}[Truncation preserves the expected-budget comparison]
\[
 V_\mu(B)\le V_{\mu,N}(B)
        \le V_\mu(B)+\rho^2B/N.
\]
\end{lemma}
\begin{proof}
The first inequality is inclusion of policy classes. For any nearly optimal
expected-budget policy, follow it until it terminates or reaches $N$ flips.
If it has not terminated by that point, output $\rho/2$. Work only decreases.
On $\{T\le N\}$ its clipped output is unchanged; on $\{T>N\}$ the squared
loss increase is at most $\rho^2$. Uniformly in $q$, Markov's inequality
gives $\mathbb P_q(T>N)\le B/N$. Take suprema and then the infimum. The
original finite expected budget implies almost-sure finite stopping work;
the truncated procedure terminates after at most $N$ observations regardless.
No deterministic work cap is inferred for the original procedure.
\end{proof}

\begin{lemma}[Uniform hidden-magnitude rounding]
For every original allocation $q$, define
$\bar q_j=\pi_j^{-1}\int_{A_j}q\,d\mu$. This allocation is admissible for
$\mu_\eta$ and preserves $\theta$ exactly. Conversely, any finite allocation
$\bar q$ lifts to $q(d)=\bar q_j$ on $A_j$, preserving feasibility and target.
For any fixed $N$-capped adaptive policy, the two transcript laws in either
comparison have total variation at most $N\eta$. Their expected work differs
by at most $N^2\eta$, and their squared risks differ by at most $\rho^2N\eta$.
\end{lemma}
\begin{proof}
On each fresh coin draw its original $D$ and $S$ jointly, and replace $D$
by its bin representative while retaining $S$. The rounded coin has exactly
the required mixture allocation $\bar q$. Its probability of heads differs
from the original by at most $\eta$, even after conditioning on both hidden
variables. Couple each queried pair of bits using a fresh common uniform.
Use the same independent policy seed. Until the first differing bit the two
policies make the same decisions, visit the same coin identities and stop
together. There are at most $N$ opportunities for a mismatch, each having
conditional probability at most $\eta$. Thus the mismatch probability, and
hence total variation, is at most $N\eta$. This holds under arbitrary
cross-coin feedback, since the latent probability difference is pointwise
bounded, not merely bounded before observations. The reverse lift uses the
same coupling. Both work and squared loss are bounded respectively by $N$
and $\rho^2$, and the target is unchanged, giving the two expectation bounds.
Free unused identities may be removed: coins are indexed by first charged
observation, so at most $N$ genuine coin tapes are needed.
\end{proof}

\begin{lemma}[A rarity-preserving finite parameter grid]
For any admissible allocation $\bar q$ under $\mu_\eta$, put
$q_j^\downarrow=\gamma\lfloor\bar q_j/\gamma\rfloor$. Then
$q^\downarrow\in Q_\gamma$,
$0\le\theta(\bar q)-\theta(q^\downarrow)\le\gamma$, and every fixed
$N$-capped policy satisfies
\[
 \TV(P_{\bar q},P_{q^\downarrow})\le N\gamma,
 \quad |\E_{\bar q}T-\E_{q^\downarrow}T|\le N^2\gamma,
\]
and
\[
 |\E_{\bar q}(Z-\theta(\bar q))^2
       -\E_{q^\downarrow}(Z-\theta(q^\downarrow))^2|
       \le\rho^2N\gamma+2\rho\gamma.
\]
\end{lemma}
\begin{proof}
Rounding downward preserves the global rarity inequality, and the target
change is at most $\sum_j\pi_j\gamma=\gamma$. On every new coin couple the
hidden type identically and the two signs by a common uniform. Their sign
disagreement probability is at most $\gamma$. If signs match, use identical
bit tapes. At most $N$ fresh coins are actually queried. Union bounding their
sign disagreements gives the transcript bound uniformly over adaptive
policies. The work bound follows from $0\le T\le N$. First hold the target
fixed to obtain the squared-loss change $\rho^2N\gamma$ from total variation;
then use $|(z-\theta)^2-(z-\theta')^2|\le2\rho|\theta-\theta'|$ for
$z,\theta,\theta'\in[0,\rho]$ to account for the target shift.
\end{proof}

\paragraph{Restoring the exact expected budget.}
If a transferred policy has uniform expected work at most $B+D$ and squared
risk at most $A$, independently run it with probability $p=B/(B+D)$ and
otherwise output $\rho/2$ without observations. This is an allowed independent
seed choice. Its work is at most $B$, and its squared risk is at most
\[
 pA+(1-p)\rho^2/4\le A+\rho^2D/[4(B+D)].                 \tag{2}
\]
This is an expected-work statement, not a high-probability or deterministic
work guarantee. The extra error is explicit.

\begin{proof}[Proof of Theorem~\ref{thm:approx}: value comparison]
Start with a nearly optimal finite-grid policy at budget $B$. For any
original $q$, average it within the bins and then round the averaged
coordinates downward. The two preceding lemmas show that running this same
policy on the original experiment has uniform squared risk at most
$G_{N,\eta,\gamma}(B)+E$ and expected work at most $B+C$. Apply (2), then
take the approximation error in the nearly optimal policy to zero. This gives
\[
 V_\mu(B)\le G_{N,\eta,\gamma}(B)+E+
                         \rho^2C/[4(B+C)].                \tag{3}
\]

For the opposite direction start with an $N$-capped original policy at
budget $B$. Every grid parameter has the bin-constant lift to the original
class. Magnitude rounding alone gives squared risk increase at most
$E_\eta=\rho^2N\eta$ and work increase at most $C_\eta=N^2\eta$.
Applying (2) now in the finite-grid experiment and then the truncation lemma
gives
\[
 G_{N,\eta,\gamma}(B)
 \le V_\mu(B)+\rho^2B/N+E_\eta+
                         \rho^2 C_\eta/[4(B+C_\eta)].      \tag{4}
\]
No grid approximation is needed in this direction: the grid is a subset of
the feasible rounded parameters. Since $C_\eta\le C$, $E_\eta\le E$, and
$D/(B+D)$ is increasing, (3)--(4) imply (1). The claimed limit follows
directly from the displayed error terms.
\end{proof}

\paragraph{Exact finite-dimensional convex formulation.}
We give the formulation to specify what a future finite numerical solve
would actually certify. A history records the ordered sequence of queried
coin identities and returned bits, with identities canonically numbered by
first observation. After $n<N$ bits its actions are stop, query any existing
coin, or query one fresh coin. At depth $N$ it must stop. This is a finite
tree, including all feedback policies. The number of nodes may be enormous;
no efficient-computation claim is made.

For each history $h$, let $\ell_q(h)$ be the product, over its distinct coins,
of the probabilities of their local bit words under the finite mixture and
grid parameter $q$. It excludes policy-choice probabilities. Let $x_h$ be a
nonnegative realization weight with $x_\varnothing=1$. At a decision node,
introduce action weights $x_{h,a}\ge0$ satisfying
\[
 \sum_a x_{h,a}=x_h.
\]
For a querying action $a$, both bit children obey
$x_{(h,a,0)}=x_{(h,a,1)}=x_{h,a}$. Chance probabilities belong in
$\ell_q$, not in these policy weights. Let $t_h=x_{h,\mathrm{stop}}$ be
the terminal weight. At depth $N$ set $t_h=x_h$. Terminal probability under
$q$ is $\ell_q(h)t_h$.

Choose scaled terminal outputs $y_h$ with $0\le y_h\le\rho t_h$. If $t_h>0$,
output $z_h=y_h/t_h$. The uniform work constraints and squared-risk objective
are the finite constraints, one pair for every $q\in Q_\gamma$,
\[
 \sum_h |h|\ell_q(h)t_h\le B,\qquad
 \sum_h\ell_q(h)\frac{(y_h-\theta(q)t_h)^2}{t_h}\le r.
\]
The sums are over terminal histories. Minimize $r$. Each quotient is the
closed quadratic perspective, assigned value zero when $t_h=y_h=0$.
The other constraints are linear, so this is a finite-dimensional convex
program; continuous terminal outputs need not be discretized.

For exactness, any such action flows produce a behavioral policy by choosing
action $a$ with probability $x_{h,a}/x_h$ when $x_h>0$; unreachable nodes may
use arbitrary decisions. Conversely, arbitrary shared and local seeds do
not create extra statistical strategies here. Conditional on a history of
actions and bits, seed consistency contributes a parameter-independent
indicator, while its observation likelihood is $\ell_q(h)$ independently
of the consistent seed. Thus the conditional law of that seed, and hence of
its next action, is independent of $q$ whenever the history is possible.
The resulting behavioral action probabilities have the same laws for all
grid parameters. At a terminal history, replace any remaining randomized
output by its conditional mean. That conditional law is again parameter
independent, and Jensen's inequality decreases squared loss for every
target. This proves equality with the stated policy minimax value, including
shared independent randomization and arbitrary cross-coin feedback.

There is always a feasible constant-output policy. All flow and scaled-output
variables range over a compact polytope; the perspectives are continuous
on $0\le y_h\le\rho t_h$ with the specified value at zero. Therefore the
minimum is attained. This existence assertion does not claim that the
program was solved in this research pass.
